\documentclass[11pt]{article}
\usepackage[margin=1in]{geometry}
\usepackage{amsmath,booktabs,hyperref}
\title{Soft Cross-Worker Occupancy Leases for Read-Only Agent Trace Queries}
\author{AssistX fleet engineering research --- working draft}
\date{October 8, 2026}
\begin{document}
\maketitle
\begin{abstract}
Global trace-history queries can perform substantial graph expansion even when page
sizes are small. A rolling-minute quota does not limit simultaneous executions.
We prototype a Redis-server-time, token-bound occupancy lease applied after
authentication and before read-only graph query execution, complementing an
earlier atomic rate guard. Offline Python tests and isolated Redis 7 Lua tests
verify a three-slot admission cap, stale-release denial and expiry-based crash
recovery on synthetic requests. A composed ten-client test with a two-request
peer quota admitted two and denied eight. The 30-second lease expires
independently of request termination; hence this is a soft concurrency guard,
not hard cancellation or proof of safe production admission.
\end{abstract}
\section{Motivation and preregistration}
Earlier isolated Neo4j 5.26 experiments over 85,000 synthetic trace groups
reported approximately $10^6$ graph DB hits for individual outcome-count and
page queries. In-flight overlap therefore matters independently of per-minute
request quotas. We preregistered a token-fenced lease design, server-owned time,
positive and negative admission tests, crash expiry, and exact-release safety
before implementing the current experiment.
\section{Method}
For active lease set $L(t)$ at server time $t$, the Redis Lua acquisition
removes entries with expiry no greater than $t$ and grants a new unique token
only when
\[
 |L(t)| < C, \quad C=3.
\]
The new entry expires at $t+30\,\mathrm{s}$. Release calls remove only the
caller-specific cryptographically random token. A distinct Redis Lua
rate-window check follows capacity acquisition. An occupied capacity
refusal cannot charge the rolling-minute rate quota. A rate refusal releases
the occupancy token before opening the graph. The request's finally block
also releases tokens after success or exception. Redis Lua scripting executes
atomically as documented by Redis \cite{redis_eval}; Redis TIME supplies
the consistency clock \cite{redis_time}.
\section{Observations}
\begin{table}[h]\centering
\caption{Disconnected Redis 7 with invented request tokens.}
\begin{tabular}{rrrr}\toprule
Simultaneous clients & Admitted & Denied & Global cap\\\midrule
1&1&0&3\\
3&3&0&3\\
5&3&2&3\\
10&3&7&3\\\bottomrule
\end{tabular}
\end{table}
In a separate composed ten-client test with shared source identity,
peer quota two and fleet minute quota five, two were admitted,
one failed the rate window and seven failed global occupancy.
Expiry and stale-release trials refused to free a successor's token.
Across the new and prior API admission/query suites, 65 Python tests passed.
These checks used only synthetic Redis tokens and mocked graph responses.
No live production database, API, Redis, or trace payload was accessed.
\section{Limitations}
This protocol is a \emph{soft} occupancy guard. A query extending past the
30-second TTL may overlap a successor without its Redis token, so the
physical number of executing queries is not bounded under arbitrary stalls.
Although each Neo4j read has a four-second timeout, persistent driver waits,
cancellation paths and host faults require independent acceptance. Rate-limit
fairness depends on a trusted reverse-proxy identity; the existing socket-peer
key may merge multiple operators. Redis failover, production authentication,
genuine browser behavior and end-to-end tool-call retention remain unproven.
The work is not deployed and does not authorize live execution or graph writes.
This arXiv-style manuscript has not been compiled, submitted or peer reviewed.
\begin{thebibliography}{9}
\bibitem{redis_eval} Redis, \emph{EVAL command}, \url{https://redis.io/docs/latest/commands/eval/}.
\bibitem{redis_time} Redis, \emph{TIME command}, \url{https://redis.io/docs/latest/commands/time/}.
\bibitem{nist92} K. Kent and M. Souppaya, \emph{Guide to Computer Security Log Management},
NIST SP 800-92 (2006). \url{https://doi.org/10.6028/NIST.SP.800-92}.
\end{thebibliography}
\end{document}
